\section{Training and Solver Integration}
\label{sec:training}

\subsection{Dataset}

All models are trained on the curated turbulence modelling dataset of \citet{McConkey2021}, which contains 902{,}601 spatial points across 38 flow cases in four geometries: square duct (14 Reynolds numbers), periodic hills (5 slope angles), converging-diverging channel, and curved backward-facing step.
Each point provides RANS fields from a $k$-$\omega$ SST \citep{Menter1994} simulation (velocity gradients, $k$, $\varepsilon$, $\omega$, $\nut$) and DNS/LES reference data (Reynolds stress anisotropy tensor $b_{ij}$).

An important characteristic of this dataset: the square duct is the only genuinely three-dimensional flow.
Its DNS reference data \citep{Pinelli2010} contains nonzero cross-plane Reynolds stresses ($\overline{u'w'}$, $\overline{v'w'}$, i.e., $b_{13}$ and $b_{23}$) that drive Prandtl's secondary flow of the second kind.
However, the RANS input fields for the duct are effectively two-dimensional: the linear eddy viscosity model ($k$-$\omega$ SST) cannot predict the secondary vortices, so the RANS cross-plane velocities $V$ and $W$ remain near zero \citep{McConkey2021}.
The remaining three geometries (hills, CDC, CBFS) have homogeneous spanwise directions and strictly two-dimensional mean flow.
The models therefore learn the 3D anisotropy correction ($b_{13}$, $b_{23}$) from the duct cases while receiving near-zero RANS cross-plane velocity gradient inputs---the correction is driven entirely by the Pope invariants and tensor basis, not by RANS-predicted secondary flow.

\subsection{Feature engineering}
\label{sec:features}

\subsubsection{Input features: 5 Pope invariants}

Following \citet{Ling2016} and \citet{Pope1975}, the inputs to all neural network models are 5 scalar invariants of the non-dimensionalized mean strain rate and rotation rate tensors.

The strain and rotation rate tensors are computed from the RANS velocity gradients:
\begin{equation}
    S_{ij} = \frac{1}{2}\left(\frac{\partial U_i}{\partial x_j} + \frac{\partial U_j}{\partial x_i}\right), \quad
    \Omega_{ij} = \frac{1}{2}\left(\frac{\partial U_i}{\partial x_j} - \frac{\partial U_j}{\partial x_i}\right),
\end{equation}
and non-dimensionalized using the turbulence time scale:
\begin{equation}
    \Shat_{ij} = \frac{k}{\varepsilon} S_{ij}, \quad \Omhat_{ij} = \frac{k}{\varepsilon} \Omega_{ij},
\end{equation}
where $k$ and $\varepsilon$ are clamped to a minimum of $10^{-30}$ to prevent division by zero.
The 5 scalar invariants are:
\begin{equation}
    \lambda_1 = \mathrm{tr}(\Shat^2), \quad
    \lambda_2 = \mathrm{tr}(\Omhat^2), \quad
    \lambda_3 = \mathrm{tr}(\Shat^3), \quad
    \lambda_4 = \mathrm{tr}(\Omhat^2 \Shat), \quad
    \lambda_5 = \mathrm{tr}(\Omhat^2 \Shat^2).
\end{equation}
These form a complete integrity basis for the independent invariants of a symmetric and an antisymmetric $3 \times 3$ tensor \citep{Pope1975}.

Input features are Z-score standardized using means and standard deviations computed from the training set only.

\subsubsection{Tensor basis}

For models predicting the full anisotropy tensor (TBNN, PI-TBNN, TBRF), the 10-tensor Pope integrity basis $T^{(n)}_{ij}$ is computed from $\Shat$ and $\Omhat$ (see Appendix~\ref{sec:appendix_basis}).
Each basis tensor is symmetric and traceless, stored as 6 independent components.

\subsubsection{Training labels}

\paragraph{TBNN / PI-TBNN / TBRF.} The full anisotropy tensor $b_{ij} = \overline{u'_i u'_j} / (2k) - \delta_{ij}/3$ from DNS, represented as 6 symmetric components.
The network predicts the raw DNS anisotropy---not the discrepancy $\Delta b = b_{\mathrm{DNS}} - b_{\mathrm{RANS}}$---following \citet{Ling2016}.

\paragraph{MLP / MLP-Large.} A scalar proxy: the Frobenius-like norm $|b| = \sqrt{\sum_c b_c^2}$ computed over the 6 stored components.

\subsection{Case-holdout split}

We adopt a case-holdout protocol where entire flow cases are held out, testing geometric generalization rather than interpolation within a single flow configuration:

\begin{itemize}
    \item \textbf{Training} (18 cases, 271{,}924 points): Square duct at 14 Reynolds numbers, periodic hills at $\alpha = \{0.5, 1.0, 1.5\}$, converging-diverging channel.
    \item \textbf{Validation} (2 cases, 23{,}967 points): Square duct $\Rey = 2000$ (interpolation), periodic hills $\alpha = 0.8$ (interpolation).
    \item \textbf{Test} (2 cases, 51{,}844 points): Periodic hills $\alpha = 1.2$ (interpolation at new angle), curved backward-facing step $\Rey = 13{,}700$ (entirely new geometry).
\end{itemize}

\subsection{Training hyperparameters}

All neural networks are trained with the Adam optimizer \citep{Kingma2015} using cosine annealing with warm restarts (initial period $T_0 = 200$ epochs, period multiplier $T_{\mathrm{mult}} = 2$), initial learning rate $10^{-3}$, minimum learning rate $10^{-6}$, and batch size 256.
Training runs for up to 1{,}000 epochs with early stopping (patience 150 epochs on validation MSE).
The best checkpoint (lowest validation MSE) is selected as the final model.
No regularization (L2 weight decay, dropout) is applied.
Full training details are provided in Appendix~\ref{sec:appendix_training}.

\subsection{PI-TBNN realizability penalty}

The PI-TBNN augments the standard MSE loss with a soft realizability penalty:
\begin{equation}
    \mathcal{L} = \mathrm{MSE}(b_{\mathrm{pred}}, b_{\mathrm{DNS}}) + \beta \cdot \frac{1}{N}\sum_{i=1}^{N} \sum_{c \in \{11,22,33\}} \left[\mathrm{ReLU}\!\left(-b_{cc}^{(i)} - \tfrac{1}{3}\right)^2 + \mathrm{ReLU}\!\left(b_{cc}^{(i)} - \tfrac{2}{3}\right)^2\right],
    \label{eq:pi_loss}
\end{equation}
which enforces the necessary realizability conditions $-1/3 \leq b_{ii} \leq 2/3$ on the diagonal components.
The penalty weight $\beta$ is linearly ramped from 0 to its target value over the first 100 epochs.
Validation loss is always computed as pure MSE (without penalty terms) for fair comparison.
We sweep $\beta \in \{0.001, 0.01\}$; results are presented in Section~\ref{sec:apriori}.

\subsection{TBRF training}

The TBRF \citep{Kaandorp2020} trains 10 independent random forests (one per basis coefficient) using scikit-learn.
Before training, the 10 tensor basis coefficients $g^{(n)}$ are extracted for each training point via the minimum-norm least-squares solution:
\begin{equation}
    \bm{g} = A^T (A A^T + \epsilon I)^{-1} \bm{b},
\end{equation}
where $A = [T^{(1)} \cdots T^{(10)}]^T$ is the $6 \times 10$ basis matrix and $\epsilon = 10^{-10}$ provides regularization.
Each forest uses 200 trees of maximum depth 20 with default scikit-learn hyperparameters and a fixed random seed of 42.
Compact variants with 1, 5, and 10 trees per coefficient are exported for solver integration experiments.

\subsection{Solver integration}
\label{sec:solver_integration}

The feature computation and inference pipeline within the solver proceeds as follows at each time step:

\begin{enumerate}
    \item \textbf{Velocity gradients}: Compute $\partial U_i / \partial x_j$ (9 components) via second-order central differences on the staggered grid---a standard stencil operation shared with EARSM.
    \item \textbf{Strain/rotation tensors}: Symmetrize and antisymmetrize the gradient tensor, non-dimensionalize by $k/\varepsilon$.
    \item \textbf{Invariants}: Compute the 5 scalar invariants via $3 \times 3$ matrix products and traces.
    \item \textbf{Normalization}: Apply the saved Z-score parameters ($\mu$, $\sigma$ from training).
    \item \textbf{Tensor basis} (TBNN/TBRF only): Compute the 10 basis tensors $T^{(n)}_{ij}$.
    \item \textbf{Inference}: Forward pass through the neural network (MLP/TBNN) or tree traversal (TBRF).
    \item \textbf{Postprocessing}: Write the predicted $\nut$ (MLP) or $\tau_{ij}$ (TBNN/TBRF) to the solver's turbulence fields.
\end{enumerate}

Steps 1--5 are implemented as GPU kernels with each cell processed independently.
For neural network inference (step 6), the computation is parallelized over (cell, neuron) pairs: each GPU thread block processes a batch of cells, with threads within a block computing different output neurons of the same layer.
Weight matrices are stored in constant memory (MLP) or read from global memory with L2 caching (TBNN).
Workspace buffers use a ping-pong pattern between layers to avoid extra memory allocation.

All sub-phases are instrumented with \texttt{TIMED\_SCOPE} profiling markers, enabling the detailed cost breakdown presented in Section~\ref{sec:cost}.
